\section{Actualización dodecafásica y memoria observable}
\label{sec:eta}
Sea $d_t$ el código entero de dirección conservado por una traza TPK.
Este código y la dirección cardinal del cursor son coordenadas distintas.
Las doce ventanas
$E_m=\{9m+1,\ldots,9m+9\}$, $0\le m<12$, dividen los 108 pasos del registro.
La orientación es
\[
\Sigma_{12}=(+,+,+,-,-,-,+,+,+,-,-,-).
\]
El evento de una ventana es la salida
\begin{equation}
a_m=\Sigma_{12}(m)
\sum_{t\in E_m}\mathbf1_{\{2,4,6,8\}}(d_t).
\label{eq:events}
\end{equation}
La notación $a_m$ evita confundir este evento con el número de Euler.
El registro fuente conserva además la orientación TRIT, la firma,
el acarreo y el segmento de historia correspondiente \cite{registro}.

\begin{proposition}[Incremento producido por el transporte]
Sea $\mathbf e_m$ la base de $\R^{12}$ y
$S\mathbf e_m=\mathbf e_{m+1\pmod{12}}$. Si
$\mathbf z_{m+1}=\mathbf z_m+a_m\mathbf e_m$ y $y_m=S^{-m}\mathbf z_m$, entonces
\begin{equation}
y_{m+1}=S^{-1}y_m+\eta_m,
\qquad
\eta_m=a_mS^{-1}\mathbf e_0.
\label{eq:eta}
\end{equation}
Después de una vuelta completa,
\begin{equation}
y_{12}=y_0+\sum_{m=0}^{11}a_m\mathbf e_m.
\label{eq:cycle}
\end{equation}
\end{proposition}
\begin{proof}
Multiplicar la actualización de $\mathbf z_m$ por $S^{-(m+1)}$ da
\[
y_{m+1}=S^{-1}y_m+a_mS^{-(m+1)}\mathbf e_m
=S^{-1}y_m+a_mS^{-1}\mathbf e_0.
\]
Al iterar, el coeficiente del evento $a_m$ es
$S^{-(12-m)}\mathbf e_0=\mathbf e_m$. Como $S^{12}=I$,
se obtiene \eqref{eq:cycle}.
\end{proof}

La negrita identifica aquí el vector acumulado
\(\mathbf z_m\in\R^{12}\), mientras que el bloque incidencial
\(z_m\) de \eqref{eq:registro-bloque-incidencial} es un entero entre
\(0\) y \(999\). El marco móvil \(y_m\) transporta el primer objeto;
la codificación decimal del segundo conserva su definición propia.

Por tanto, $\eta_m$ no es una fuerza libre elegida para imponer una
conservación: es la expresión vectorial del evento TPK. La igualdad de fase
después de una vuelta no cancela los eventos acumulados. Este lector
conserva los recuentos, mientras el registro completo conserva además su
orden y procedencia. Se distingue $\eta_m$, incremento vectorial, de
$\eta_{\rm ret}$, mantisa escalar de acción, y de $\eta_{\rm el}$,
parámetro de forma que se introducirá en la sección~\ref{sec:ellipse}.

\subsection{Balance cuadrático}
Definimos
\[
\nabla_3=I-S^3,\quad \nabla_4=I-S^4,\quad
B=\nabla_3^{\mathsf T}\nabla_3+\nabla_4^{\mathsf T}\nabla_4,\qquad
\mathcal M(y)=\tfrac12 y^{\mathsf T}By.
\]
Los lectores orientados del registro firmado cumplen
$(D_jz)_m=z_m-z_{m+j}$; con la convención
$S\mathbf e_m=\mathbf e_{m+1}$, se tiene
$D_j=I-S^{-j}=\nabla_j^{\mathsf T}$. Ambos sentidos producen la misma
forma cuadrática $B$, pues las potencias de $S$ conmutan.
Se cumplen $S^{\mathsf T}BS=B$ y $B_{00}=4$. Sustituyendo
\eqref{eq:eta} se deduce exactamente
\begin{equation}
\mathcal M(y_{m+1})-\mathcal M(y_m)
=a_m(By_m)_0+2a_m^2.
\label{ii:mem:balance}
\end{equation}
Asimismo, la carga $q(y)=\sum_jy_j$ satisface
$q(y_{m+1})-q(y_m)=a_m$.
La prueba es la expansión del polinomio cuadrático tras retirar la
permutación $S^{-1}$. Ningún valor metrológico entra en ese cálculo.
La forma $\mathcal M$ mide este registro; su interpretación como una
energía física determinada requeriría el lector dimensional correspondiente.

\subsection{Resolventes del transporte y conservación de la memoria}
\label{subsec:ii-memoria-resolvente}

La historia prolongada determina las ventanas
\(I_m=\{9m+1,\ldots,9m+9\}\), para \(m\geq0\).
Con la numeración de esta sección, \(I_m=E_m\) durante el primer ciclo.
Sea \(\sigma_m\in\{-1,1\}\) la orientación conservada por esa historia;
para \(0\leq m<12\), coincide con \(\Sigma_{12}(m)\). Se define
\[
 a_m=\sigma_m\sum_{t\in I_m}\mathbf1_{\{2,4,6,8\}}(d_t),
 \qquad |a_m|\leq9.
\]
El registro y el marco móvil se prolongan por
\[
 \mathbf z_{m+1}=\mathbf z_m+a_m\mathbf e_{m\bmod12},\qquad y_m=S^{-m}\mathbf z_m.
\]
El estado inicial \(\mathbf z_0\) conserva la historia previa; vale cero para
un prefijo inicialmente vacío. En este apartado escribimos
\(R=S^{-1}\) para el transporte del marco. La actualización
\eqref{eq:eta} se mantiene en toda ventana. La orientación y los eventos
posteriores se obtienen de la historia prolongada: el retorno del marco
no presupone periodicidad de la secuencia de eventos.

\begin{proposition}[Identidad de truncación con estado terminal]
Para $N\geq1$, sean
\[
Y_N(u)=\sum_{m=0}^N y_mu^m,\qquad
H_{\eta,N}(u)=\sum_{m=0}^{N-1}\eta_mu^m.
\]
Se cumple la identidad polinómica exacta
\begin{equation}
(I-uR)Y_N(u)=y_0+uH_{\eta,N}(u)-u^{N+1}Ry_N.
\label{mem:identidad-finita}
\end{equation}
\end{proposition}
\begin{proof}
En grados $1\leq m\leq N$, el coeficiente del miembro izquierdo es
$y_m-Ry_{m-1}=\eta_{m-1}$, por \eqref{eq:eta}.
Sus coeficientes en grados cero y $N+1$ son $y_0$ y $-Ry_N$,
respectivamente. La comparación de coeficientes prueba la igualdad.
\end{proof}

El término terminal conserva la memoria transportada al cortar una historia.
La identidad permite comparar una evaluación finita con su prolongación
manteniendo explícito el estado que queda en la frontera.

\subsubsection{Correspondencia con el transporte del núcleo común}

La numeración de ventanas de la sección~\ref{subsec:memoria-resolvente}
y la empleada aquí describen la misma partición de la historia. La
correspondencia es
\[
 \underbrace{I_m=E_{m+1}}_{\text{índice del núcleo común}}
 \quad\longleftrightarrow\quad
 \underbrace{I_m=E_m}_{\text{índice de esta sección}},
 \qquad 0\leq m<12.
\]
El subíndice cambia el nombre de la ventana; sus nueve transiciones,
su orientación \(\sigma_m=\Sigma_{12}(m)\), su evento \(a_m\) y
su posición en el registro permanecen iguales. Para toda prolongación,
\(R=S^{-1}\), \(y_m=S^{-m}\mathbf z_m\) y la
actualización~\eqref{eq:eta} coinciden con
\eqref{mem:actualizacion}. La orientación y la secuencia de eventos
conservan la dependencia respecto de la historia, también después del
retorno del marco.

Esta identificación permite aplicar directamente la
proposición~\ref{mem:prop-serie}: la serie~\eqref{mem:serie}, la
cota~\eqref{mem:cota} y la convergencia uniforme de sus derivadas
corresponden a este mismo registro y al mismo reloj de ventanas. La
identidad finita~\eqref{mem:identidad-finita} determina además el término
terminal al truncar esa serie. Sus hipótesis mantienen la historia inicial
y el evento acotado \(|a_m|\leq9\).

Al reutilizar el registro, la proposición~\ref{mem:prop-schur} conserva
conjuntamente el operador, la fuente y el lector mediante
\eqref{mem:schur-fuente} y~\eqref{mem:schur-lector}. El ejemplo
\eqref{mem:ejemplo-ciclo} precisa el efecto de las excursiones omitidas
por una compresión. La composición afín~\eqref{mem:cociclo-tres}, los
bloques intermedios~\eqref{mem:schur-intermedio} y la identidad
\eqref{mem:schur-transitivo} garantizan la coherencia al concatenar
segmentos o eliminar varias capas. Sus demostraciones completas residen
en la sección~\ref{subsec:memoria-resolvente}, con los dominios, las
fuentes y las aplicaciones de lectura allí especificados.

El balance~\eqref{ii:mem:balance} conserva la correspondencia entre los
lectores orientados del registro firmado y la forma cuadrática común.
La distinción entre el vector acumulado \(\mathbf z_m\) y el bloque
decimal incidencial \(z_m\) sigue siendo la establecida en esta sección.
Las identidades de transporte se aplican a la memoria así tipada; no
determinan por sí solas coeficientes analíticos adicionales.


\section{Sello dodecafásico y lectura racional reversible}
\label{sec:sello}

El registro terminal tiene una coordenada visible y una coordenada
firmada. El sello se obtiene de esta última mediante la transformación
de Hadamard por órbitas; una lectura escalar posterior permite recuperar
el sello cuando se conserva su marco de representación. Esta propiedad
precisa qué significa aquí una definición estructural de una constante:
se explicitan los mapas de producción y recuperación, además de su valor.

Sean $U\in\Z^{12}$ la coordenada firmada producida por el registro,
$H_4$ la matriz de Hadamard desarrollada en la sección~\ref{sec:action}
y $P_{\rm orb}$ la permutación que agrupa
\[
 (1,4,7,10),\qquad(2,5,8,11),\qquad(3,6,9,12).
\]
Defínase
\[
 \mathcal H_{12}=P_{\rm orb}^{-1}
 (H_4\oplus H_4\oplus H_4)P_{\rm orb}.
\]
Se cumple $\mathcal H_{12}^2=4I$. Cuando
$\mathcal H_{12}U\equiv0\pmod4$, el sello integral es
\[
 K=\tfrac14\mathcal H_{12}U,
 \qquad U=\mathcal H_{12}K.
\]
Las diferencias transversales $D_3K,D_4K$ y la carga $Q(K)$ recuperan
los canales especificados por el registro \cite{registro}.
El sello es así una coordenada reconstruible de ese registro orientado;
la profundidad y la historia completa siguen siendo campos distintos.

\begin{proposition}[Recuperación de bloques y transporte del origen]
Fijadas una base entera $b\ge2$ y doce posiciones ordenadas, sea
$K_j\in\{0,\ldots,b-1\}$. Defínanse
\[
 I(K)=\sum_{j=1}^{12}K_jb^{12-j},\qquad
 \kappa(K)=\frac{I(K)}{b^{12}-1}.
\]
La lectura exacta $\kappa$ determina unívocamente los doce bloques.
Para la rotación $\rho K=(K_2,\ldots,K_{12},K_1)$,
\[
 \boxed{\kappa(\rho K)=b\kappa(K)-K_1.}
\]
\end{proposition}
\begin{proof}
El entero $I=(b^{12}-1)\kappa$ se descompone de manera única en doce
posiciones de base $b$, rellenando con ceros las posiciones iniciales
necesarias. Por otra parte,
\[
 I(\rho K)=bI(K)-K_1(b^{12}-1).
\]
Dividir por $b^{12}-1$ da la ley de transformación. Para el bloque
constante $K_j=b-1$, la lectura es $\kappa=1$; el tipo de representación
periódica fijado conserva igualmente su recuperación.
\end{proof}

En la base $1000$ del registro, los bloques de tres cifras permanecen
recuperables, incluidos los ceros iniciales. La lectura racional
periódica conserva más que una aproximación decimal truncada.
Al modificar el origen de fase, cambia generalmente su valor según
la identidad anterior; al completar doce rotaciones, retorna.
El avance de memoria del estado TPK es una transformación adicional:
la periodicidad del lector no identifica las historias acumuladas.

\begin{apunte}{lectura reversible e información}
La recuperación exacta de los bloques es una propiedad del lector
racional en su marco ordenado. La complejidad de palabras, la entropía
de una dinámica simbólica, la entropía del estado reducido y la
producción termodinámica son objetos distintos. Una relación entre
ellos requiere especificar la distribución o el estado y el mapa de
transporte correspondiente; la reversibilidad del sello, por sí sola,
no identifica esos funcionales.
\end{apunte}

